\begin{textsample}{POS Dim 3 – ro – Score 22.00 – ro/ro\_inf\_35.txt}  \label{ex:f3_pos_139}
Espaços, tempos, quantidades, relações e transformações Desde cedo as \textbf{crianças} vivem inseridas em espaços e tempos de diferentes dimensões, em um mundo constituído de fenômenos naturais e socioculturais.
Desde muito \textbf{pequenas} participam de momentos na interação com a família que envolvem conhecimentos relacionados a quantidades, tempos, espaços, relações e transformações, quando, por exemplo, respondem perguntas referentes à sua idade, observam o uso do dinheiro quando acompanham os \textbf{adultos} às compras, são pesadas e medidas ; calculam a distância entre um local e outro, exploram objetos, mudam os canais da televisão, \textbf{contam} seus \textbf{brinquedos}, \textbf{contam} seus calçados, assim como elaboram conhecimentos sobre os números, espaços, formas e medidas, etc., vivenciados no seu cotidiano.
Ainda procuram se situar em diversos espaços ( rua, bairro, cidade etc. ) e tempos ( \textbf{dia} e noite, hoje, ontem e amanhã etc. ). \textbf{Demonstram} também \textbf{curiosidade} sobre o mundo físico ( seu próprio corpo, os fenômenos atmosféricos, os animais, as plantas, as transformações da natureza, os diferentes tipos de materiais e as possibilidades de sua manipulação etc. ) e o mundo sociocultural ( as relações de parentesco e sociais entre as pessoas que conhece ; como vivem e em que trabalham essas pessoas ; quais suas tradições e seus costumes ; a diversidade entre elas etc. ), isso vai propiciando à \textbf{criança} o raciocínio lógico, chegando à descoberta de novos saberes.
Conforme as experiências oportunizadas, nas quais as \textbf{crianças} têm oportunidade de explorar diferentes características e propriedades de objetos, materiais, \textbf{brinquedos}, jogos de construção, ao que se refere à forma, tamanho, espessura etc., explorando, \textbf{manipulando}, observando, \textbf{contando} e medindo os objetos, elas lidam com noções de quantidades, classes, medidas e formas, ampliando suas habilidades.
Experiências de apreciar uma \textbf{pintura}, desenhar, localizar -se, deslocar -se, ler, escrever, \textbf{brincar} e muitas outras ampliam as noções da \textbf{criança} de espaço.
Além disso, nessas experiências e em muitas outras, as \textbf{crianças} também se \textbf{deparam}, frequentemente, com conhecimentos matemáticos ( contagem, ordenação, relações entre quantidades, dimensões, medidas, comparação de \textbf{pesos} e de comprimentos, avaliação de distâncias, reconhecimento de formas geométricas, conhecimento e reconhecimento de numerais cardinais e ordinais etc. ) que igualmente aguçam a \textbf{curiosidade}.
Para isso, é importante que a intencionalidade educativa para a primeira \textbf{infância} venha como forma de encantamento e de maneira a aguçar a \textbf{curiosidade} das \textbf{crianças} para elaborarem suas hipóteses, questionar, discordarem, compararem, em busca de respostas para suas investigações.
Na integralidade dos campos de experiências as \textbf{crianças} podem vivenciar momentos de apreciar uma \textbf{pintura}, desenhar, localizar -se, ler, escrever, \textbf{brincar}, e muitas outras situações que ampliam as noções da \textbf{criança} de espaço.
A \textbf{criança}, por meio das experiências, vai se apropriando e percebendo o mundo de números nas observações cotidianas que faz dos endereços, das placas de \textbf{carros} e motos, documentos, preços, número de calçados e roupas, situações onde é necessário \textbf{contar} e/ou recitar numericamente, registrar, de forma não convencional, pontos em um jogo, entre outros.
As situações de grandezas e medidas podem ser oportunizadas por meio de medidas dos espaços, de uma mesa para confeccionar uma toalha, nas receitas \textbf{culinárias} que podem envolver as \textbf{crianças} na produção de alguma alimentação, quando poderão comparar medidas utilizando diferentes objetos, bem como poderão observar as relações de quantidade, \textbf{volume}, densidade e aspectos físicos dos produtos.
Ao observar os espaços e os objetos e as muitas relações que estes possuem com as atribuições matemáticas, as \textbf{crianças} vão inferindo as relações e transformações no meio físico.
Portanto, promover experiências na Educação \textbf{Infantil} nas quais as \textbf{crianças} falam, descrevem, narram, explicam, pesquisam e exploram, torna -se, assim, requisito fundamental para a construção e ampliação de saberes.
Por exemplo, as vivências cotidianas das \textbf{crianças} ao construir um castelo como cenário de um espaço ambientado, procurar um sapo, borboleta no jardim, \textbf{cuidar} de plantas e de animais, colecionar objetos, além de fortalecer sua autonomia, podem ser ricas oportunidades para a construção do raciocínio lógico, de noções de tempo e espaço, de classificações, etc., para a \textbf{percepção} de mudanças e transformações nos objetos e materiais observados ou \textbf{manuseados}, e para o desenvolvimento da imaginação pela \textbf{criança}.
Por isso, as \textbf{instituições} precisam promover experiências nas quais as \textbf{crianças} possam ter oportunidades de ampliar seus conhecimentos do mundo físico, natural e sociocultural e possam utilizá -los em seu cotidiano.

% matched lemmas: adulto, brincar, brinquedo, carro, contar, criança, cuidar, culinário, curiosidade, demonstrar, deparar, dia, infantil, infância, instituição, manipular, manusear, pequeno, percepção, peso, pintura, volume
\end{textsample}
